\IfFileExists{stacks-project.cls}{%
\documentclass{stacks-project}
}{%
\documentclass{amsart}
}

% For dealing with references we use the comment environment
\usepackage{verbatim}
\newenvironment{reference}{\comment}{\endcomment}
%\newenvironment{reference}{}{}
\newenvironment{slogan}{\comment}{\endcomment}
\newenvironment{history}{\comment}{\endcomment}

% For commutative diagrams we use Xy-pic
\usepackage[all]{xy}

% We use 2cell for 2-commutative diagrams.
\xyoption{2cell}
\UseAllTwocells

% We use multicol for the list of chapters between chapters
\usepackage{multicol}

% This is generally recommended for better output
\usepackage{lmodern}
\usepackage[T1]{fontenc}

% For cross-file-references
\usepackage{xr-hyper}

% Package for hypertext links:
\usepackage{hyperref}

% For any local file, say "hello.tex" you want to link to please
% use \externaldocument[hello-]{hello}
\externaldocument[introduction-]{introduction}
\externaldocument[conventions-]{conventions}
\externaldocument[sets-]{sets}
\externaldocument[categories-]{categories}
\externaldocument[topology-]{topology}
\externaldocument[sheaves-]{sheaves}
\externaldocument[sites-]{sites}
\externaldocument[stacks-]{stacks}
\externaldocument[fields-]{fields}
\externaldocument[algebra-]{algebra}
\externaldocument[brauer-]{brauer}
\externaldocument[homology-]{homology}
\externaldocument[derived-]{derived}
\externaldocument[simplicial-]{simplicial}
\externaldocument[more-algebra-]{more-algebra}
\externaldocument[smoothing-]{smoothing}
\externaldocument[modules-]{modules}
\externaldocument[sites-modules-]{sites-modules}
\externaldocument[injectives-]{injectives}
\externaldocument[cohomology-]{cohomology}
\externaldocument[sites-cohomology-]{sites-cohomology}
\externaldocument[dga-]{dga}
\externaldocument[dpa-]{dpa}
\externaldocument[sdga-]{sdga}
\externaldocument[hypercovering-]{hypercovering}
\externaldocument[schemes-]{schemes}
\externaldocument[constructions-]{constructions}
\externaldocument[properties-]{properties}
\externaldocument[morphisms-]{morphisms}
\externaldocument[coherent-]{coherent}
\externaldocument[divisors-]{divisors}
\externaldocument[limits-]{limits}
\externaldocument[varieties-]{varieties}
\externaldocument[topologies-]{topologies}
\externaldocument[descent-]{descent}
\externaldocument[perfect-]{perfect}
\externaldocument[more-morphisms-]{more-morphisms}
\externaldocument[flat-]{flat}
\externaldocument[groupoids-]{groupoids}
\externaldocument[more-groupoids-]{more-groupoids}
\externaldocument[etale-]{etale}
\externaldocument[chow-]{chow}
\externaldocument[intersection-]{intersection}
\externaldocument[pic-]{pic}
\externaldocument[weil-]{weil}
\externaldocument[adequate-]{adequate}
\externaldocument[dualizing-]{dualizing}
\externaldocument[duality-]{duality}
\externaldocument[discriminant-]{discriminant}
\externaldocument[derham-]{derham}
\externaldocument[local-cohomology-]{local-cohomology}
\externaldocument[algebraization-]{algebraization}
\externaldocument[curves-]{curves}
\externaldocument[resolve-]{resolve}
\externaldocument[models-]{models}
\externaldocument[functors-]{functors}
\externaldocument[equiv-]{equiv}
\externaldocument[pione-]{pione}
\externaldocument[etale-cohomology-]{etale-cohomology}
\externaldocument[proetale-]{proetale}
\externaldocument[relative-cycles-]{relative-cycles}
\externaldocument[more-etale-]{more-etale}
\externaldocument[trace-]{trace}
\externaldocument[crystalline-]{crystalline}
\externaldocument[spaces-]{spaces}
\externaldocument[spaces-properties-]{spaces-properties}
\externaldocument[spaces-morphisms-]{spaces-morphisms}
\externaldocument[decent-spaces-]{decent-spaces}
\externaldocument[spaces-cohomology-]{spaces-cohomology}
\externaldocument[spaces-limits-]{spaces-limits}
\externaldocument[spaces-divisors-]{spaces-divisors}
\externaldocument[spaces-over-fields-]{spaces-over-fields}
\externaldocument[spaces-topologies-]{spaces-topologies}
\externaldocument[spaces-descent-]{spaces-descent}
\externaldocument[spaces-perfect-]{spaces-perfect}
\externaldocument[spaces-more-morphisms-]{spaces-more-morphisms}
\externaldocument[spaces-flat-]{spaces-flat}
\externaldocument[spaces-groupoids-]{spaces-groupoids}
\externaldocument[spaces-more-groupoids-]{spaces-more-groupoids}
\externaldocument[bootstrap-]{bootstrap}
\externaldocument[spaces-pushouts-]{spaces-pushouts}
\externaldocument[spaces-chow-]{spaces-chow}
\externaldocument[groupoids-quotients-]{groupoids-quotients}
\externaldocument[spaces-more-cohomology-]{spaces-more-cohomology}
\externaldocument[spaces-simplicial-]{spaces-simplicial}
\externaldocument[spaces-duality-]{spaces-duality}
\externaldocument[formal-spaces-]{formal-spaces}
\externaldocument[restricted-]{restricted}
\externaldocument[spaces-resolve-]{spaces-resolve}
\externaldocument[formal-defos-]{formal-defos}
\externaldocument[defos-]{defos}
\externaldocument[cotangent-]{cotangent}
\externaldocument[examples-defos-]{examples-defos}
\externaldocument[algebraic-]{algebraic}
\externaldocument[examples-stacks-]{examples-stacks}
\externaldocument[stacks-sheaves-]{stacks-sheaves}
\externaldocument[criteria-]{criteria}
\externaldocument[artin-]{artin}
\externaldocument[quot-]{quot}
\externaldocument[stacks-properties-]{stacks-properties}
\externaldocument[stacks-morphisms-]{stacks-morphisms}
\externaldocument[stacks-limits-]{stacks-limits}
\externaldocument[stacks-cohomology-]{stacks-cohomology}
\externaldocument[stacks-perfect-]{stacks-perfect}
\externaldocument[stacks-introduction-]{stacks-introduction}
\externaldocument[stacks-more-morphisms-]{stacks-more-morphisms}
\externaldocument[stacks-geometry-]{stacks-geometry}
\externaldocument[moduli-]{moduli}
\externaldocument[moduli-curves-]{moduli-curves}
\externaldocument[examples-]{examples}
\externaldocument[exercises-]{exercises}
\externaldocument[guide-]{guide}
\externaldocument[desirables-]{desirables}
\externaldocument[coding-]{coding}
\externaldocument[obsolete-]{obsolete}
\externaldocument[fdl-]{fdl}
\externaldocument[index-]{index}

% Theorem environments.
%
\theoremstyle{plain}
\newtheorem{theorem}[subsection]{Theorem}
\newtheorem{proposition}[subsection]{Proposition}
\newtheorem{lemma}[subsection]{Lemma}

\theoremstyle{definition}
\newtheorem{definition}[subsection]{Definition}
\newtheorem{example}[subsection]{Example}
\newtheorem{exercise}[subsection]{Exercise}
\newtheorem{situation}[subsection]{Situation}

\theoremstyle{remark}
\newtheorem{remark}[subsection]{Remark}
\newtheorem{remarks}[subsection]{Remarks}

\numberwithin{equation}{subsection}

% Macros
%
\def\lim{\mathop{\mathrm{lim}}\nolimits}
\def\colim{\mathop{\mathrm{colim}}\nolimits}
\def\Spec{\mathop{\mathrm{Spec}}}
\def\Hom{\mathop{\mathrm{Hom}}\nolimits}
\def\Ext{\mathop{\mathrm{Ext}}\nolimits}
\def\SheafHom{\mathop{\mathcal{H}\!\mathit{om}}\nolimits}
\def\SheafExt{\mathop{\mathcal{E}\!\mathit{xt}}\nolimits}
\def\Sch{\mathit{Sch}}
\def\Mor{\mathop{\mathrm{Mor}}\nolimits}
\def\Ob{\mathop{\mathrm{Ob}}\nolimits}
\def\Sh{\mathop{\mathit{Sh}}\nolimits}
\def\NL{\mathop{N\!L}\nolimits}
\def\CH{\mathop{\mathrm{CH}}\nolimits}
\def\proetale{{pro\text{-}\acute{e}tale}}
\def\etale{{\acute{e}tale}}
\def\QCoh{\mathit{QCoh}}
\def\Ker{\mathop{\mathrm{Ker}}}
\def\Im{\mathop{\mathrm{Im}}}
\def\Coker{\mathop{\mathrm{Coker}}}
\def\Coim{\mathop{\mathrm{Coim}}}

% Boxtimes
%
\DeclareMathSymbol{\boxtimes}{\mathbin}{AMSa}{"02}

%
% Macros for moduli stacks/spaces
%
\def\QCohstack{\mathcal{QC}\!\mathit{oh}}
\def\Cohstack{\mathcal{C}\!\mathit{oh}}
\def\Spacesstack{\mathcal{S}\!\mathit{paces}}
\def\Quotfunctor{\mathrm{Quot}}
\def\Hilbfunctor{\mathrm{Hilb}}
\def\Curvesstack{\mathcal{C}\!\mathit{urves}}
\def\Polarizedstack{\mathcal{P}\!\mathit{olarized}}
\def\Complexesstack{\mathcal{C}\!\mathit{omplexes}}
% \Pic is the operator that assigns to X its picard group, usage \Pic(X)
% \Picardstack_{X/B} denotes the Picard stack of X over B
% \Picardfunctor_{X/B} denotes the Picard functor of X over B
\def\Pic{\mathop{\mathrm{Pic}}\nolimits}
\def\Picardstack{\mathcal{P}\!\mathit{ic}}
\def\Picardfunctor{\mathrm{Pic}}
\def\Deformationcategory{\mathcal{D}\!\mathit{ef}}
\usepackage{cleveref}
\crefname{originaldefinition}{original definition}{original definitions}
\Crefname{originaldefinition}{Original definition}{Original definitions}
\makeatletter
\newlabel{situation-complexes}{{0DLE}{source}{}{sourcepointer-1}{}}
\newlabel{situation-complexes@cref}{{[originaldefinition][0][]0DLE}{[1][0][]source}{}{}{}}
\makeatother
\crefname{originalremark}{original remark}{original remarks}
\Crefname{originalremark}{Original remark}{Original remarks}
\makeatletter
\newlabel{artin-remark-strong-effectiveness}{{0CXT}{source}{}{sourcepointer-2}{}}
\newlabel{artin-remark-strong-effectiveness@cref}{{[originalremark][0][]0CXT}{[1][0][]source}{}{}{}}
\makeatother
\crefname{originallemma}{original lemma}{original lemmas}
\Crefname{originallemma}{Original lemma}{Original lemmas}
\makeatletter
\newlabel{lemma-complexes}{{0DLD}{source}{}{sourcepointer-3}{}}
\newlabel{lemma-complexes@cref}{{[originallemma][0][]0DLD}{[1][0][]source}{}{}{}}
\makeatother
\crefname{originaltheorem}{original theorem}{original theorems}
\Crefname{originaltheorem}{Original theorem}{Original theorems}
\makeatletter
\newlabel{spaces-flat-theorem-existence-derived}{{0DIQ}{source}{}{sourcepointer-4}{}}
\newlabel{spaces-flat-theorem-existence-derived@cref}{{[originaltheorem][0][]0DIQ}{[1][0][]source}{}{}{}}
\makeatother
\crefname{originallemma}{original lemma}{original lemmas}
\Crefname{originallemma}{Original lemma}{Original lemmas}
\makeatletter
\newlabel{more-algebra-lemma-limit-henselian}{{0CT7}{source}{}{sourcepointer-5}{}}
\newlabel{more-algebra-lemma-limit-henselian@cref}{{[originallemma][0][]0CT7}{[1][0][]source}{}{}{}}
\makeatother
\crefname{originallemma}{original lemma}{original lemmas}
\Crefname{originallemma}{Original lemma}{Original lemmas}
\makeatletter
\newlabel{spaces-more-morphisms-lemma-henselian-relatively-perfect}{{0DL5}{source}{}{sourcepointer-6}{}}
\newlabel{spaces-more-morphisms-lemma-henselian-relatively-perfect@cref}{{[originallemma][0][]0DL5}{[1][0][]source}{}{}{}}
\makeatother
\crefname{originallemma}{original lemma}{original lemmas}
\Crefname{originallemma}{Original lemma}{Original lemmas}
\makeatletter
\newlabel{lemma-complexes-open-neg-exts-vanishing}{{0DLC}{source}{}{sourcepointer-7}{}}
\newlabel{lemma-complexes-open-neg-exts-vanishing@cref}{{[originallemma][0][]0DLC}{[1][0][]source}{}{}{}}
\makeatother
\title{Strong formal effectiveness for the stack of relatively perfect complexes}
\author{The Stacks Project Authors}
\date{stacks-project:quot.tex:local-snapshot-20261012; GNU FDL 1.2 or later}
\begin{document}
\maketitle
\setcounter{subsection}{1}
\noindent\textit{Editorial note.} All delivered mathematical text is the author's own, copied from the pinned Stacks Project chapter file quot.tex; no statement or proof step was rewritten, reordered or removed. The counted span is the author environment and its own proof environment. References inside the span are handed to the assembler as navigation stubs that carry the canonical tag of the cited tag only; whether each of them is load-bearing is recorded on the stub itself.\par\medskip
\noindent\textit{Allowed standard prerequisites (editorial).}\begin{itemize}
\item Situation situation-complexes (canonical tag 0DLE). cited at statement line 5847. Standing situation: B = S locally Noetherian.
\item Remark artin-remark-strong-effectiveness (canonical tag 0CXT). cited at statement line 5849. The property being established: strong formal effectiveness.
\item Lemma lemma-complexes (canonical tag 0DLD). cited at proof lines 5866 and 5884. Condition (2) for the objects and for the limit.
\item Theorem spaces-flat-theorem-existence-derived (canonical tag 0DIQ). cited at proof line 5870. Existence of E with the prescribed derived pullbacks.
\item Lemma more-algebra-lemma-limit-henselian (canonical tag 0CT7). cited at proof line 5877. (R, Ker(R -> R\_1)) is a henselian pair.
\item Lemma spaces-more-morphisms-lemma-henselian-relatively-perfect (canonical tag 0DL5). cited at proof lines 5879-5880. E is R-perfect.
\item Lemma lemma-complexes-open-neg-exts-vanishing (canonical tag 0DLC). cited at proof line 5885. The negative-Ext vanishing locus is open.
\end{itemize}\medskip
\ifdefined\hypertarget\hypertarget{sourcepointer-1}{}\fi
\noindent\textit{Source reference.} Original definition 0DLE: sources/stacks-chow/tags.txt, source lines 17044--17044. Cited at statement line 5847: the lemma is phrased in this situation; load-bearing.\par\medskip
\ifdefined\hypertarget\hypertarget{sourcepointer-2}{}\fi
\noindent\textit{Source reference.} Original remark 0CXT: sources/stacks-chow/tags.txt, source lines 16218--16218. Cited at statement line 5849: the property being established; load-bearing.\par\medskip
\ifdefined\hypertarget\hypertarget{sourcepointer-3}{}\fi
\noindent\textit{Source reference.} Original lemma 0DLD: sources/stacks-chow/tags.txt, source lines 17043--17043. Cited at proof lines 5866 and 5884: condition (2) for the objects and for the limit; load-bearing twice.\par\medskip
\ifdefined\hypertarget\hypertarget{sourcepointer-4}{}\fi
\noindent\textit{Source reference.} Original theorem 0DIQ: sources/stacks-chow/tags.txt, source lines 16950--16950. Cited at proof line 5870: the object E of D(O\_\{X\_R\}) with E\_n its derived pullbacks; load-bearing.\par\medskip
\ifdefined\hypertarget\hypertarget{sourcepointer-5}{}\fi
\noindent\textit{Source reference.} Original lemma 0CT7: sources/stacks-chow/tags.txt, source lines 16057--16057. Cited at proof line 5877: the henselian pair, giving Ker in the Jacobson radical; load-bearing.\par\medskip
\ifdefined\hypertarget\hypertarget{sourcepointer-6}{}\fi
\noindent\textit{Source reference.} Original lemma 0DL5: sources/stacks-chow/tags.txt, source lines 17035--17035. Cited at proof lines 5879-5880: E is R-perfect; load-bearing.\par\medskip
\ifdefined\hypertarget\hypertarget{sourcepointer-7}{}\fi
\noindent\textit{Source reference.} Original lemma 0DLC: sources/stacks-chow/tags.txt, source lines 17042--17042. Cited at proof line 5885: the negative-Ext vanishing locus is open; load-bearing.\par\medskip
\begin{lemma}
\label{lemma-complexes-strong-effectiveness}
In Situation \ref{situation-complexes} assume $B = S$ is locally Noetherian.
Then strong formal effectiveness in the sense of
Artin's Axioms, Remark \ref{artin-remark-strong-effectiveness}
holds for $p : \Complexesstack_{X/S} \to (\Sch/S)_{fppf}$.
\end{lemma}

\begin{proof}
Let $(R_n)$ be an inverse system of $S$-algebras with surjective transition
maps whose kernels are locally nilpotent. Set $R = \lim R_n$.
Let $(\xi_n)$ be a system of objects of $\Complexesstack_{X/B}$
lying over $(\Spec(R_n))$. We have to show $(\xi_n)$ is effective, i.e.,
there exists an object $\xi$ of $\Complexesstack_{X/B}$ lying over
$\Spec(R)$.

\medskip\noindent
Write $X_R = \Spec(R) \times_S X$ and $X_n = \Spec(R_n) \times_S X$.
Of course $X_n$ is the base change of $X_R$ by $R \to R_n$. Since $S = B$,
we see that $\xi_n$ corresponds simply to an $R_n$-perfect object
$E_n \in D(\mathcal{O}_{X_n})$ satisfying condition (2) of
Lemma \ref{lemma-complexes}. In particular $E_n$ is pseudo-coherent.
The isomorphisms $\xi_{n + 1}|_{\Spec(R_n)} \cong \xi_n$
correspond to isomorphisms $L(X_n \to X_{n + 1})^*E_{n + 1} \to E_n$.
Therefore by
Flatness on Spaces, Theorem \ref{spaces-flat-theorem-existence-derived}
we find a pseudo-coherent object $E$ of $D(\mathcal{O}_{X_R})$
with $E_n$ equal to the derived pullback of $E$ for all $n$
compatible with the transition isomorphisms.

\medskip\noindent
Observe that $(R, \Ker(R \to R_1))$ is a henselian pair, see
More on Algebra, Lemma \ref{more-algebra-lemma-limit-henselian}.
In particular, $\Ker(R \to R_1)$ is contained in the Jacobson radical of $R$.
Then we may apply More on Morphisms of Spaces, Lemma
\ref{spaces-more-morphisms-lemma-henselian-relatively-perfect}
to see that $E$ is $R$-perfect.

\medskip\noindent
Finally, we have to check condition (2) of Lemma \ref{lemma-complexes}.
By Lemma \ref{lemma-complexes-open-neg-exts-vanishing}
the set of points $t$ of $\Spec(R)$ where the negative self-exts of $E_t$
vanish is an open. Since this condition is true in $V(\Ker(R \to R_1))$
and since $\Ker(R \to R_1)$ is contained in the Jacobson radical of $R$
we conclude it holds for all points.
\end{proof}

\end{document}
